

\section{SE-PIM Programming Model}
\label{sec:prog-model}
%\hjc{What does this achieve??? This is like a instruction manual for PIM programming. No explanation on its purpose, advantages, and our contribution. Unnecessary long and pointless. REWRITE}

In this section, we explain the programming model of \thename. Our proposed framework presents an abstraction on top of the \thename architecture we introduced. Our programming model supports the two design objectives of \thename that we outline in \cref{subsec:objectives}: a secure memory extension and a secure large data accelerator for the secure enclave. We elaborate the software components of \thename that help an enclave interact with \thename in \cref{subsec:software}. Then, we go through the design of our programming model that enable our design objectives in \cref{subsec:prog-model}.

%We first describe the interfaces between the host enclave and \thename in \cref{subsec:interfaces}. Then, we illustrate the programming model of offloading computations to \thename and using \thename as secure memory extension through code examples (\cref{subsec:prog-model}).



\subsection{Software Components}
\label{subsec:software}
%The software components of \thename includes the facilitates the usage of \thename from a secure enclave. 
As the programming model of SGX enclaves does not support a trusted I/O path, the communication channel between the enclave and the devices has to be set up by the untrusted kernel and host application. Hence, there are untrusted and trusted components in our software infrastructure. Our trusted software includes the enclave host program, the trusted \thename user library, and the \thename kernel. Our untrusted components are the kernel driver and the untrusted user library. The untrusted components only set up the secure communication channel; after a secure session is established, all communication between the enclave and \thename is encrypted and authenticated.
%\subsection{Overview}
%\hdr{\thename host application.}

\hdr{\thename driver and user library.}
The communication interfaces between the host and the \thename unit are initialized by the driver and the user library. The driver maps the physical addresses corresponding to the memory banks in the \thename-enabled memory modules into the host kernel address space. It also maps the registers of each \thename unit to expose the device control channel to the host. In order to provide access to \thename, the driver creates a virtual file interface to the host library in the user space (e.g., \texttt{/dev/pim0}). A user library at the host side is required to allow an in-host program to communicate and send workloads to \thename. The \thename library has two components, a trusted part that resides in the enclave code and an untrusted part in the host program. The untrusted part of the library initializes the communication interfaces by obtaining the device driver's memory banks and control registers. The trusted library contains functions that facilitate the communication between the enclave and \thename units (e.g., transfer data to the memory bank and sending commands to \thename).



%The driver and the user library help the host enclave establish communication channels with the \thename units.  
\hdr{Enclave host program and \thename kernel.}
The enclave host program contains the confidential computation logic that uses \thename as its accelerator. \thename kernel includes protected functions to be executed inside \thename. The developer of the host enclave must select the computation (e.g., functions) to be offloaded, then implement and cross-compile them for \thename platform. 
% The kernel objects are linked with a runtime library with helper functions for PIM operations (e.g., sending DMA requests to the memory bank). 

The similar to the interfaces of \baseline, host enclave communicates with the \thename kernel running on the \thename unit core through two physical interfaces (described in \cref{subsec:pim-baseline}): \emph{memory-mapped registers} and \emph{direct memory access (DMA)}. On the two interfaces, we establish two communication \emph{channels}. The \emph{command channel} (CMD CH) use the memory-mapped registers to issue \emph{commands} (listed in \autoref{tab:commands}) to the PIM core to control its execution.  

The \emph{parameters channel} (PARAM CH in \autoref{fig:interfaces}) uses DMA to send parameters of the commands. It functions as a secure communication channel between the host enclave and a \thename unit. Messages through the channel are encrypted using a shared session key (\texttt{SESSION\_KEY}) yielded during the initial attestation and key exchange~(\cref{subsec:attestation}). On \texttt{EXECUTE} commands, the parameters become the arguments for execution that is defined by the \thename kernel.


%It uses the memory mappings obtained by the untrusted driver and user library to send commands and parameters to \thename.

%A common header file is also employed for shared data structures (e.g., \texttt{SE\_PIM\_params\_t}) between \thename kernels and the host process.

%\subsection{Using \thename as a Secure Memory Extension}
% \hjc{Too many subsections A through L is too many.}







%Before the execution of a \thename kernel, the host sends parameters containing the function name (FN) to be executed and the arguments (FN ARGS) to the \emph{parameter channel} (PARAM CH). 


% We also demonstrate an overview of our execution model in \autoref{fig:interfaces}. 
% 
% \hj{The previous claim was that zero-copy is possible?}

%\hj{You said the regions are accessed just like normal memory. But we use DMA?}

% \todo{Mention Zero-copy...}


% \hjc{This requires explanation of RA components in the hardware}



% \hjc{Maybe merge secure comm ch. with remote attestation?}

% \hdr{The encrypted parameter channel.}
% \hdr{Secure communication channel.} 
%The host enclave communicates with the PIM kernel running on the \thename unit core through two physical interfaces: \emph{memory-mapped registers} and \emph{direct memory access (DMA)}. % For DMA, the host process maps the buffer into its own memory space. 
% On the two interfaces, we establish two communication \emph{channels}. The \emph{command channel} use the memory-mapped registers to issue \emph{commands} (listed in \autoref{tab:commands}) to the \thename core to control its execution. For each \thename core, there is a corresponding control register. For instance, writing an \texttt{EXECUTE} command to a control register tells a \thename core to start executing the kernel. 


%\hdr{Data loading through the data channel.}  Therefore, we allow directly writing of data into the memory bank before starting PIM computation for efficient data transfer between host and PIM. After the initial transfer finishes, the memory bank can enter a \emph{lockdown} mode to protect data confidentiality and integrity (\cref{subsec:access-control}).


%\subsection{Extending Enclave Memory with \thename}

\subsection{SE-PIM Usage Model} 
\label{subsec:prog-model}
\label{subsec:accelerator-model}

% SE_PIM_cert_t *cert = se_pim.get_cert(); |\label{line:attest_start}|
% remote_attest_and_exchange_keys(se_pim, cert);
% se_pim.offload_kernel("kernel.elf.enc"); |\label{line:attest_end}|
%transfer_data_to_mem(ptr_A, ptr_B);
% // Use the secure interface to access data from memory
% se_pim.read(data_read, ptr_A, DATA_SZ);
% se_pim.memcpy(ptr_A, ptr_B, DATA_SZ);
% void *ptr_B = se_pim.transfer(data_B, size_B);



% \begin{figure}[t]
    \begin{listing}[h]
    \centering
\begin{minted}[linenos=true,escapeinside=!!]{c}
// Initialize and and attest a SE-PIM unit
SE_PIM *se_pim = SE_PIM_init(MEM_MODULE_ID,
                                       BANK_ID, 
                                       "kernel.elf.enc"); !\label{line:object}!
// Transfer the data to SE-PIM memory bank
void *ptr = se_pim.transfer(obj, DATA_SZ); !\label{line:transfer_start}!

se_pim.lock_mem(); !\label{line:transfer_end}!

// Send command, execute and get result
SE_PIM_params_t params = {.fn = "increment", 
                                   .data_ptr = ptr, 
                                   .size = DATA_SZ};

for (int i = 0; i < ITERATIONS; i++)
    SE_PIM_result_t *result = se_pim.execute(params);

se_pim.read(result, ptr, DATA_SZ); !\label{line:secure_access}!
\end{minted}
    \caption{Host-side code snippet illustrating computation offloading to \thename.},
    \label{lst:accelerator},
\end{listing}

% \end{figure}

\autoref{lst:accelerator} demonstrates the code example of a host-side enclave that uses a \thename unit as an accelerator. The enclave first obtains an object representing a \thename unit with the library function call \texttt{SE\_PIMInit} (line \ref{line:object}). It performs remote attestation and key exchange during initialization, then offloads the encrypted \thename kernel to the \thename unit. Then, it transfers encrypted data object into the memory banks through the direct mapping and instructs the \thename unit to \emph{lock} the memory bank (lines \ref{line:transfer_start}-\ref{line:transfer_end}). The enclave then generates a parameter structure (\texttt{SE\_PIM\_params\_t}) consisting of the in-kernel function name, the pointer to the data to be processed, and the size. It then commands the \thename unit to start execution in multiple iterations while keeping the data protected inside the memory. Finally, the host read the result from the memory bank (line \ref{line:secure_access}).


%\hdr{Zero-copy}
\hdr{Extending enclave memory with \thename.} 
To support our objective as a secure memory extension for processor enclaves, our user library provides functions to offload data into the protection of \thename and access data without leaking data access patterns on the memory bus. Our data loading function (line \ref{line:transfer_start} of \autoref{lst:accelerator}) uses direct memory mappings to efficiently load data into the \thename memory banks without compromising security. As we discussed in \cref{subsec:access-control}, the act of sequential data loading does not generate distinguishable access patterns. Furthermore, we observe in \cref{subsec:secure-access-eval} that there are overheads when using the controlled memory access channel that is side-channel resistant, which would inevitably create overheads for large data transfers. After the initial transfer, the enclave commands the memory bank to enter \emph{lockdown} to protect the data inside memory (explained in \cref{subsec:access-control}).

After DRAM lockdown, data within memory is accessed through the user library functions \texttt{read()}, \texttt{write()} and \texttt{memcpy()}. The functions internally create encrypted memory requests to the controlled memory access channel (demonstrated in \cref{subsec:secure-access}), and let \thename perform the actual data access. Hence, the data access pattern is protected from any observing attacker. In the example in \autoref{lst:accelerator}, the host enclave use the function \texttt{read()} to read the incremented data block from memory without leaking the requested address and the type of operation (line \ref{line:secure_access}).

\begin{figure}[h]
    \centering
    \begin{minted}[mathescape,escapeinside=||]{c}
void increment(data_ptr, size){
    char local_mem[BATCH_SZ];
    SE_PIM_result_t result;
    result.count = 0;
    // Request data from bank memory 
    for (int i = 0; i < size / BATCH_SZ; i++){
        sepim_dma(data_ptr + BATCH_SZ * i, local_mem, |\label{line:dma_read}|
                      BATCH_SZ, BANK_READ_DATA);
        // Process data in local memory
        for(int j = 0; j < BATCH_SZ; j++){ |\label{line:process_start}|
            local_mem[i] += 1; 
            result.count++;} |\label{line:process_end}|
        sepim_dma(local_mem, data_ptr + BATCH_SZ * i, |\label{line:dma_write}|
                     BATCH_SZ, BANK_WRITE_DATA);
    }
    return_result(&result);|\label{line:return}|
}
\end{minted}
    \caption{\thename kernel code snippet showing the kernel fetches data into local memory, increments, writes updated data to bank memory, and finally returns the results to the enclave.}
    \label{lst:kernel}
\end{figure}

\hdr{Offloading confidential computation to \thename.}
\autoref{lst:kernel} demonstrates the kernel logic used in \autoref{lst:accelerator}. 
After remote attestation and key exchange, the encrypted \thename kernel binary is offloaded through the secure channel as the argument to the command \texttt{LOAD\_KERNEL}. The kernel is loaded into \thename local memory and used by the PIM core for execution. Then, the PIM core takes the measurement (i.e., the hash) of PIM local memory and sends the measurement to the host enclave. The host enclave can use this measurement to verify the software state of a \thename unit. After loading the PIM kernel, the \thename unit goes into an idle state and waits for future commands (e.g., \texttt{EXECUTE} or \texttt{LOCK\_MEM}). 

\thename makes use of the \emph{zero-copy} model of PIM-based computation, as described in \cref{subsec:pim-baseline}. Data is not moved for computation, but only the \emph{pointers} to the data are embedded in function arguments. \thename uses the pointers to fetch the actual data from the bank memory into its local memory and perform the requested operation on the data. This eliminates the data transfers for each computation and makes use of the fast internal bus of PIM systems~\cite{hmc,primer-on-pim}. Moreover, the computation does not leave any pattern on the memory bus. This is only possible on PIM-based computation since CPUs, or any accelerators on the peripheral bus, inevitably access memory with patterns unless using ORAM primitives. To command a \thename unit to execute the kernel, the host sends the encrypted parameters through the encrypted parameter channel; then sends an \texttt{EXECUTE} command to the corresponding \thename unit. After execution, \thename kernel encrypts and returns the execution results to the enclave through the same encrypted channel (\ref{line:return}).

%There are two types of computational results from PIM. 
%First, the \thename kernel might update the data inside the memory due to the computation. For instance, the kernel in ~\autoref{lst:kernel} increments the values in an in-memory array by one. The host retrieves the data with the controlled memory channel proposed in \cref{subsec:secure-access} to hide the address side-channels. Second, the computation results can be information obtained from the data inside the memory content. This type of result is similar to the return value of functions calls. For instance, in \autoref{lst:kernel}, the kernel calculates a statistic (the count of elements) that is obtained based on the data. The result is encrypted with the session key and written to the secure communication channel. 
%\hdr{Loading \thename kernels.} 

%The kernel fetches data from the memory bank to its local memory, increments it, and then writes it back to the same location. Data is processed in batches, for each batch, the kernel initiates DMA transfers through built-in functions \texttt{dma\_request()} to load and decrypt data from the memory bank into \thename local memory (line \autoref{line:dma-read}). It then processes data in the local memory (lines \autoref{line:process-start}-\autoref{line:process-end}). Finally, the kernel updates the memory bank's content with a DMA transfer in the opposite direction (line \autoref{line:dma-write}).
%The kernel also keeps track of the number of elements it accessed. It returns the count of elements as the computation results through the secure channel (line \autoref{line:return}).



%\hj{leave only absolutely necessary code}



% the index of the targeted memory bank (housing a \thename unit) and memory module 

%The programming model of \thename also adapted by many previous works~\cite{upmem,tesseract,impica,neural-network-training}. 







%\hdr{Initialization and \thename kernel loading.} 
%The untrusted host library first obtains the memory mapping of the memory bank and the control registers from the driver with a memory mapping request (\texttt{mmap()}) on the virtual file interface of the driver. Then, the trusted runtime library initializes an object (line \autoref{line:object} in \autoref{lst:accelerator}) that contains a \thename unit mappings, the memory allocation status, and functions that help the enclave communicate with \thename units. When multiple \thename units are employed, a \thename object is created for each \thename unit. 
%\hdr{\thename attestation and keys exchange.}
%The host enclave then uses the communication interfaces (i.e., the virtual mapping) to send attestation requests to \thename units. After, the host enclave exchanges a session key and a data decryption key with the \thename unit. \thename and the host enclave use the session key to encrypt all messages on the secure communication channel from this point onward. We describe the attestation and key exchange processes in more detail in \cref{subsec:attestation}.  %Since \thename units use the same root endorsement key

%\hdr{Transfering encrypted data to \thename memory banks.} Different from the programming model of GPUs that requires a DMA buffer to transfer data from and to the VRAM~\cite{graviton}, the host use uncached memory mapping to transfer encrypted data directly to the \thename memory banks. After data transfer finished, the host enclave commands \thename to enable DRAM lockdown (\cref{subsec:access-control}).  %The data transfer only occurs before the execution of the \thename kernel; after, the access to the memory bank is \emph{locked}, except for accesses to the DMA-based communication channel.

%\hdr{Execution of \thename kernels.}  %Securing the execution of the kernels are discussed in \cref{sec:secure-acc}. 
%Note that no additional data transfer is needed for the execution of the kernel. The host only sends the \emph{pointer} of data to be processed inside memory, which the kernel will use to fetch the data by itself with DMA transfers. 


%\hdr{Execution of \thename kernel and results retrieval.} 







% Note that all of the commands on \thename share the same parameter passing method described in the next section. However, the
% and key storage that only allows indirect interactions with the root endorsement key (\texttt{EK}) through a predefined set of commands.

